\section{恢复路径与有向账本}\label{sec:environment}

尽管形式结果适用于任何一维目标相对调整，全文的贯穿解释仍是相对于分阶段治疗目标的成瘾存量。形式路径是一份恢复日记；可达性问题询问：给定的阈值跨越与端点是否能够共同属于同一条连续历史。

\subsection{路径与模块}

令
\[
 X=[\underline x,\overline x]\subset\R,
 \qquad \underline x<0<\overline x,
\]
并固定目标为 $0$。记 $R_L=-\underline x$、$R_R=\overline x$。所谓\emph{有限步路径}，是连续映射 $\gamma:[0,1]\to X$，使得存在 $[0,1]$ 的某个有限分割，并且 $\gamma$ 在每个分割单元上单调。允许存在常值片段。

四个通行模块为
\[
 \M=\{T_L,T_R,A_L,A_R\}.
\]
符号 $T$ 表示朝向目标的运动，$A$ 表示远离目标的运动；下标表示所在一侧。对 $m\in\M$，令 $R_m$ 表示相应一侧上的最大目标距离。

对 $0\le a<b\le R_m$，令 $p_m(a,b)$ 表示跨越距离区间 $[a,b]$ 的典范单调通行。趋向目标 通行从距离 $b$ 移向 $a$；远离目标 通行从 $a$ 移向 $b$。端点约定是不对称的：
\begin{itemize}
 \item 趋向目标 通行 $p_{T_s}(0,b)$ 记录目标达成，因此覆盖 $[0,b]$；
 \item 远离目标 通行 $p_{A_s}(0,b)$ 记录经 $(0,b]$ 离开目标的过程，并且在 $0$ 处不另计单点通行。
\end{itemize}

\subsection{通行剖面}

每个单调片段在其所属模块中贡献一个阶梯函数。对所有片段求和得到有向通行剖面
\[
 h^m_\gamma:[0,R_m]\to\N_0,
 \qquad m\in\M.
\]
对于穿孔区间 $(a,b]$，$h^m_\gamma(r)$ 计数模块 $m$ 中有多少片段跨越状态水平 $r$。在目标处，$h^{T_s}_\gamma(0)$ 是从侧 $s$ 到达目标的次数，而 $h^{A_s}_\gamma(0)=0$。具体账本为
\[
 \ell(\gamma)
 =\bigl(h^{T_L}_\gamma,h^{T_R}_\gamma,h^{A_L}_\gamma,h^{A_R}_\gamma\bigr)
 \in\Led.
\]
形式上，$\Led$ 是四个整数阶梯函数群的直积。趋向目标 一侧的群由 $\1_{[0,b]}$ 与满足 $0<a<b$ 的 $\1_{(a,b]}$ 生成，而 远离目标 一侧的群由满足 $0\le a<b$ 的 $\1_{(a,b]}$ 生成。实际账本构成该群中非负且可达的子集。
账户不记录各片段的具体先后顺序，也不记录遍历这些片段所花费的时间，但保留重复次数。

一个有限账本在其断点所出现的状态上诱导出一个有限的带标签有向多重图。每一次对相邻状态区间的单位跨越都对应一条按所在一侧和方向标记的弧。

\begin{proposition}[记录通行账户的可行性]\label{prop:reachability}
令 $L$ 为有限的四模块阶梯账本，并令 $s,t\in X$ 为给定的有序端点。
\begin{enumerate}
 \item 若 $L\ne0$，则 $L$ 能由一条从 $s$ 到 $t$ 的有限步路径生成，当且仅当其非孤立的带标签多重图连通，并且
 \[
  \deg^+(v)-\deg^-(v)
  =\1_{\{v=s\}}-\1_{\{v=t\}}
 \]
 对每个顶点 $v$ 都成立；此外，当 $s=t$ 时，共同端点 $s$ 必须属于非孤立支撑。
 \item 若 $L=0$，则存在一条从 $s$ 到 $t$ 的路径生成 $L$，当且仅当 $s=t$；此时实现路径在弱时间变换意义下都是常值路径。
\end{enumerate}
对于满足第~1 项条件的非零账本，带标签多重图的一条 Euler 迹给出一条实现路径。
\end{proposition}

\begin{proof}[证明思路]
一条实现路径必须按照记录的重复次数遍历每条带标签弧，因此强制要求非孤立支撑连通并满足所显示的端点不平衡条件。反过来，把重复次数替换成相应数量的平行弧以后，这些条件恰好就是有向 Euler 迹存在的判据。沿 Euler 迹读取顶点并在相邻状态之间插值得到一条实现用的有限步路径；零账本只能由常值路径实现。完整证明见附录~\ref{app:universal}。
\end{proof}

命题~\ref{prop:reachability} 是基本的定义域限制。账本的形式和未必满足其连通性或端点平衡条件，因此实际路径域并不对加法封闭。

\subsection{偏好}

原始偏好 $\succeq$ 是所有有限步路径集合 $\Pth$ 上的条件弱序，其中非通行后果保持固定。状态 $x$ 处的常值路径记为 $c_x$。除非某个结果明确把注意力限制在有限菜单上，否则 $\Pth$ 表示 $X$ 上完整的有限步路径域。特别地，它包含所有常值路径、反向路径、端点相容的拼接、典范通行、典范目标外往返、触及目标的循环，以及把任意典范区间细分成任意给定有限个非空典范区间所得的路径。各项假设承担四种不同角色：弱序和 BPCC 约束实际比较的可加一致性；PVBC 和 PROC 分别约束实数完备化与穿孔连续性；典范区间丰富性是定义域条件；方向单调性则给目标相对模块提供经济学方向。

\begin{assumption}[面向目标的方向单调性]\label{ass:directional}
对任意非空典范区间 $I$，每一个趋向目标 通行 $p_{T_s}(I)$ 都严格优于目标处的常值路径 $c_0$，而 $c_0$ 又严格优于每一个远离目标 通行 $p_{A_s}(I)$。在 BPCC 下，所有常值路径具有相同的零通行账户，因此彼此无差异。
\end{assumption}

方向单调性赋予预先给定的目标一个行为方向：在维持的背景后果条件下，局部朝向目标的运动优于停留，而局部远离目标的运动更差。在正则表示中，它推出非负的目标达成残差；但达成残差是否严格为正仍是偏好的额外性质。BPCC 则从相同账本推出纤维内无差异。
